\addchapnonumber{A.\ \ Chimie troposphérique de polluants clés}\label{annex:chimieAtmo}


\vspace*{\fill}
\begin{singlespace}
\textbf{Contenu de l'annexe}

    
    
~

\startcontents[chapters]
\printcontents[chapters]{}{1}{}
\end{singlespace} 
\vspace*{\fill}

\newpage


%%%%%%%%%%%%%%%%%%%%%%%%%%%%%%%%%%%%%%%%%%%%%%

Dans cette annexe, nous revenons sur les bases de la chimie troposphérique des polluants clés émis par les feux. Ce récapitulatif se veut général, mais permet d'appréhender les mécanismes en jeu lors de l'émission de polluants par les feux. 
Les éléments de cette partie reposent sur les livres suivants : \citet{jacobIntroductionAtmosphericChemistry1999,delmasPhysiqueChimieLatmosphere2005,brasseurModelingAtmosphericChemistry2017,akimotoHandbookAirQuality2023}.



%%%%%%%%%%%%%%
\section{Monoxyde de carbone}
%%%%%%%%%%%%%%%%%

Le monoxyde de carbone est, d'une part, émis directement lors de la combustion incomplète de biomasse (voir partie~\ref{sec:impact_feux_polluants}) et d'hydrocarbures fossiles (environ 2/3 du CO atmosphérique), et est, d'autre part, le produit de l'oxydation dans l'atmosphère du méthane ou d'autres hydrocarbures (environ 1/3 du CO atmosphérique). 

Le principal puits du CO est son oxydation par le radical hydroxyle (OH) : 

\begin{chemequations}
\begin{align}
\ce{CO + OH  ->[O_2] CO_2 + HO_2 }
\label{equ:R1}	
\end{align}
\end{chemequations}

La réaction du CO avec l'OH représente 40\,\% de l'élimination de l'OH dans la troposphère \citep{lelieveldGlobalTroposphericHydroxyl2016}. La concentration de CO influence donc fortement la capacité oxydante de l'atmosphère terrestre. Du fait de cette réaction le CO affecte la durée de vie et l'abondance du méthane et d'autres composés. La durée de vie du CO dans la troposphère est en moyenne de 38 jours sous les tropiques, de 65 jours dans les régions extra-tropicales de l'hémisphère Nord et de 86 jours dans celle de l'hémisphère Sud \citep{lelieveldGlobalTroposphericHydroxyl2016}. Sa durée de vie est suffisamment longue pour permettre son transport à l'échelle intercontinentale, ce qui en fait un traceur utile pour le transport à longue distance des panaches de combustion.

Le CO est toxique à des concentrations très élevées qui sont de l'ordre de 500 fois les concentrations de fond. Il n'est donc pas directement un gaz problématique pour la qualité de l'air. Cependant, l'oxydation du CO en présence d'oxydes d'azote contribue à la formation d'ozone, un polluant très important pour la qualité de l'air (voir partie~\ref{sec:Qair_general}). 



\stopcontents[chapters]
